
\section{SUBPART C -- ADDITIONAL PROTECTIONS FOR SUBJECTS IN PHYSICAL AI CLINICAL INVESTIGATIONS}

This Subpart C establishes additional protections for human subjects in clinical investigations that involve Physical AI systems. The prior 21 CFR Part 50 reserved Subpart C. This adaptation utilizes Subpart C to address protections specific to Physical AI systems that do not fit naturally into the existing Subpart A (General Provisions) or Subpart B (Informed Consent) sections. These protections are derived from the physical-ai-oncology-trials repository (v2.3.0, DOI: 10.5281/zenodo.18445179), the national MCP-PAI oncology trials standard (v1.2.0, DOI: 10.5281/zenodo.18869776), the ICH E6(R3) Physical AI adaptation (DOI: 10.5281/zenodo.18973368), and the federated learning framework (DOI: 10.5281/zenodo.18840880).

\subsection{\S~50.30 Physical AI System Safety Requirements}

\subsubsection{Pre-Procedure Safety Matrix}

(a) Before any Physical AI system may begin a clinical procedure on a human subject, a pre-procedure safety matrix must be completed and documented. The safety matrix consists of the following verification checks, all of which must pass before the procedure may proceed:

\begin{description}[leftmargin=1.0cm,style=sameline]
\item[(1)] \textit{Authorization.} A valid authorization token must be issued by the trialmcp-authz server, with scope appropriate to the planned procedure type, the assigned robot agent, and the specific subject. The token must not be expired and must not have been revoked.
\item[(2)] \textit{Patient Identity Verification.} The subject's identity must be verified against the clinical record accessed through the trialmcp-fhir server. The verification must confirm that the subject enrolled in the investigation is the same individual presenting for the procedure.
\item[(3)] \textit{Clinical Data Availability.} The subject's relevant clinical data must be accessible through the trialmcp-fhir server, including current diagnosis, treatment history, allergies, relevant laboratory results, and any contraindications for robotic procedures. Data must be de-identified per HIPAA Safe Harbor before delivery to the Physical AI system.
\item[(4)] \textit{Imaging Data Access.} Where the procedure requires imaging guidance, the subject's relevant imaging data must be accessible through the trialmcp-dicom server with appropriate role-based modality restrictions (CT, MR, PT as applicable). Imaging metadata must be normalized and patient identifiers hashed per the DICOM integration specifications.
\item[(5)] \textit{Robot Readiness and USL Threshold.} The Physical AI system must be verified as operationally ready, with current calibration, software version matching the protocol specification, and all safety systems functional. The platform's current USL score must meet or exceed the minimum threshold for the procedure type: surgical robots require USL 7.0 or above, therapeutic positioning and diagnostic needle placement require USL 6.0 or above, rehabilitative exoskeletons require USL 4.0 or above, and companion monitoring systems require USL 3.0 or above.
\item[(6)] \textit{Environmental Checks.} The clinical environment must meet the Physical AI system's operational requirements, including adequate space, appropriate lighting, electromagnetic compatibility, network connectivity for MCP server communication, and absence of interference sources that could affect robotic control or sensor performance.
\item[(7)] \textit{Simulation Validation.} The Physical AI system must have undergone simulation validation for the planned procedure type, with documented cross-framework behavioral consistency demonstrating trajectory deviations below 2mm and force discrepancies below 0.5N across at least two recognized simulation frameworks.
\item[(8)] \textit{Digital Twin Synchronization.} Where a digital twin model is used for procedure planning or guidance, the twin must be synchronized with the subject's most recent clinical data, calibrated to acceptable accuracy thresholds, and verified as computationally current before the procedure begins.
\end{description}

\subsubsection{Runtime Safety Monitoring}

(b) During the execution of any clinical procedure involving a Physical AI system, the following runtime safety monitoring requirements must be maintained continuously:

\begin{description}[leftmargin=1.0cm,style=sameline]
\item[(1)] \textit{Continuous Authorization.} The Physical AI system's authorization token must remain valid throughout the procedure. The trialmcp-authz server must be queried at regular intervals (not exceeding 60 seconds) to verify that the token has not been revoked and that no policy changes have occurred that would affect the system's permissions.
\item[(2)] \textit{Real-Time Audit Recording.} Every significant action taken by the Physical AI system during the procedure must be recorded in the hash-chained audit trail via the trialmcp-ledger server. Audit records must include the action type, timestamp, actor identity, resource affected, outcome, and a SHA-256 hash linking the record to the previous entry in the chain.
\item[(3)] \textit{Emergency Stop Availability.} An emergency stop capability must be immediately accessible to the supervising clinical personnel at all times during the procedure. Activation of the emergency stop must immediately halt all Physical AI system motion, preserve the current state for post-incident analysis, record an emergency stop audit event, and transition the task order to a failed state.
\item[(4)] \textit{Force and Torque Limit Monitoring.} For Physical AI systems that make physical contact with the subject, continuous monitoring of applied forces and torques must be maintained. The monitoring system must detect any exceedance of the pre-specified safety limits and must trigger an immediate protective response (speed reduction, controlled stop, or emergency stop) when limits are approached or exceeded.
\item[(5)] \textit{Human Oversight.} A qualified clinical professional with authority to intervene, override, or terminate the Physical AI system's operations must be physically present and attentive throughout the procedure. The supervising professional must have received documented training on the specific Physical AI system and must have access to override controls.
\end{description}

\subsubsection{Post-Procedure Requirements}

(c) Following the completion (or termination) of any clinical procedure involving a Physical AI system, the following post-procedure requirements must be fulfilled:

\begin{description}[leftmargin=1.0cm,style=sameline]
\item[(1)] \textit{Complete Audit Trail.} The complete hash-chained audit trail for the procedure must be finalized and verified for chain integrity via the trialmcp-ledger server. The final audit record must document the procedure outcome (completed, cancelled, or failed), the total procedure duration, and any safety events that occurred.
\item[(2)] \textit{Provenance Chain.} A complete data provenance record must be created via the trialmcp-provenance server, documenting the directed acyclic graph (DAG) of data lineage from the subject's clinical data through all processing steps, procedure execution, and outcome recording. Each node in the provenance DAG must include a SHA-256 fingerprint for integrity verification.
\item[(3)] \textit{Outcome Recording.} The clinical outcome of the procedure must be recorded in a structured format, linked to the audit trail and provenance chain, and made available for monitoring and safety review.
\item[(4)] \textit{Safety Incident Reporting.} Any safety events that occurred during the procedure (emergency stop activations, force limit exceedances, software anomalies, unexpected AI behaviors, equipment malfunctions) must be documented in a structured safety incident report, linked to the relevant audit records, and reported to the sponsor and IRB in accordance with the investigation's safety reporting procedures.
\item[(5)] \textit{USL Score Update.} The procedure outcome must be incorporated into the ongoing USL assessment for the Physical AI platform, contributing to the clinical trial collaboration dimension (Dimension 4) of the USL score. Aggregate procedure outcomes across the investigation inform the platform's evolving readiness profile.
\end{description}

\subsubsection{Task-Order Lifecycle}

(d) Every clinical procedure involving a Physical AI system must be governed by a task order that progresses through the following defined lifecycle states:

\begin{description}[leftmargin=1.0cm,style=sameline]
\item[(1)] \textit{Scheduled.} The procedure has been scheduled and a task order created, specifying the procedure type, assigned robot platform, subject identifier (pseudonymized), scheduled time, and required safety checks.
\item[(2)] \textit{Safety Check.} The pre-procedure safety matrix specified in paragraph (a) is being executed. The task order remains in this state until all required safety checks pass or until a check failure terminates the sequence.
\item[(3)] \textit{In Progress.} All safety checks have passed and the Physical AI system is actively executing the clinical procedure. Runtime safety monitoring specified in paragraph (b) is active throughout this state.
\item[(4)] \textit{Completed.} The procedure has been successfully completed and all post-procedure requirements specified in paragraph (c) have been fulfilled.
\item[(5)] \textit{Cancelled.} The procedure was cancelled before execution began (e.g., due to safety check failure, subject withdrawal, or scheduling change).
\item[(6)] \textit{Failed.} The procedure was terminated during execution due to a safety event, emergency stop activation, equipment malfunction, or clinical decision to abort. All post-procedure requirements must still be fulfilled for failed procedures.
\end{description}

\subsubsection{Forbidden Operations}

(e) The following operations are strictly prohibited for Physical AI systems operating in clinical investigations:

\begin{description}[leftmargin=1.0cm,style=sameline]
\item[(1)] No Physical AI system may begin a clinical procedure without completing the pre-procedure safety matrix specified in paragraph (a).
\item[(2)] No Physical AI system may operate on a subject without a valid, non-expired, non-revoked authorization token from the trialmcp-authz server.
\item[(3)] No Physical AI system may skip or disable audit recording during any clinical procedure. Every procedure must produce a complete, hash-chained audit trail.
\item[(4)] No Physical AI system may operate in a clinical investigation if its current USL score falls below the minimum threshold specified for the procedure type.
\item[(5)] No Physical AI system may directly modify clinical systems (EHR, PACS, laboratory information systems) without explicit authorization mediated through the appropriate MCP server. All clinical system interactions must be audited and traceable.
\item[(6)] No Physical AI system may continue a clinical procedure after an emergency stop has been activated without completing a new full safety check cycle as specified in paragraph (a). The original task order must be marked as failed, and a new task order must be created for any resumption of the procedure.
\end{description}

\subsection{\S~50.31 IRB Review of Physical AI Investigations}

(a) In addition to the IRB review requirements specified elsewhere in this chapter, IRBs reviewing clinical investigations that involve Physical AI systems must evaluate the following:

\begin{description}[leftmargin=1.0cm,style=sameline]
\item[(1)] \textit{Physical AI System Safety Profiles.} The IRB must review the safety profile of each Physical AI system specified in the protocol, including mechanical safety data, software reliability evidence, AI model validation results, simulation testing outcomes, and any prior clinical experience with the system.
\item[(2)] \textit{USL Scores.} The IRB must review the USL readiness score for each robotic platform and verify that the score meets or exceeds the minimum threshold for the proposed procedure type. The IRB must assess whether the four-dimension USL evaluation (simulation switching, AI integration, cross-robot sharing, clinical trial collaboration) adequately characterizes the platform's readiness for the planned clinical use.
\item[(3)] \textit{Simulation Validation Reports.} The IRB must review simulation validation evidence demonstrating that the Physical AI system has been tested in at least two recognized simulation frameworks with documented cross-framework behavioral consistency within acceptable tolerances (trajectory deviation below 2mm, force discrepancy below 0.5N).
\item[(4)] \textit{Human Oversight Plans.} The IRB must evaluate the adequacy of the human oversight plan for the investigation, including the qualifications of supervising personnel, the availability and accessibility of override controls, the emergency stop provisions, and the procedures for managing Physical AI system failures during clinical procedures.
\item[(5)] \textit{Informed Consent Adequacy.} The IRB must assess whether the informed consent document adequately addresses all Physical AI-specific elements required under \S\S~50.20 and 50.25 of this part, including the disclosure of robot types, USL scores, simulation validation, digital twin usage, MCP data protections, and the subject's right to non-Physical AI alternatives.
\item[(6)] \textit{Robot Capability Profiles.} The IRB must review the machine-readable robot capability profile for each Physical AI platform, as defined in the national MCP-PAI standard, and verify that the platform's documented capabilities, safety prerequisites, and required MCP tools are appropriate for the proposed clinical procedures.
\end{description}

(b) IRBs must have access to expertise in Physical AI systems sufficient to evaluate protocols involving robotic systems, AI/ML models, digital twins, and simulation technologies. This expertise may be provided through IRB membership, consultation with qualified experts, or collaboration with institutions that have established Physical AI review capabilities.

(c) The IRB site policy template from the national MCP-PAI standard provides a reference framework for developing institution-specific IRB review procedures for Physical AI investigations. Institutions should adapt this template to their local governance requirements.

\subsection{\S~50.32 Ongoing Consent and Subject Notification}

(a) Subjects enrolled in clinical investigations involving Physical AI systems must be notified of significant changes to the Physical AI systems used in their care. Significant changes include, but are not limited to:

\begin{description}[leftmargin=1.0cm,style=sameline]
\item[(1)] Software updates to robotic control systems or AI models that alter the system's behavior, capabilities, or safety characteristics.
\item[(2)] AI model version changes where the updated model was trained on different data, uses a modified architecture, or produces materially different outputs than the version disclosed in the original consent.
\item[(3)] Changes to the platform's USL score that move the score below the minimum threshold for the procedure type or that change the platform's readiness band classification.
\item[(4)] Substitution of a different robotic platform than the one originally disclosed in the informed consent document, even if the substitute platform is of the same type and meets the USL threshold.
\item[(5)] Changes to the MCP server configuration, security policies, or data protection mechanisms that affect how the subject's data is processed, stored, or shared.
\end{description}

(b) When significant changes are identified, the investigator must notify the IRB and determine whether re-consent is required. Changes governed by a Predetermined Change Control Plan (PCCP) that have been pre-approved by the IRB may be implemented with notification rather than re-consent, provided the PCCP was disclosed in the original consent document and the subject consented to the PCCP-governed update process.

(c) Subjects must be informed of their right to withdraw consent for Physical AI system interactions at any time, including after learning of a significant change. Withdrawal of consent for Physical AI interactions does not require withdrawal from the overall investigation if non-Physical AI alternatives are available.

\subsection{\S~50.33 Data Protection for Physical AI Investigations}

(a) \textit{De-Identification.} All clinical data accessed by Physical AI systems must be de-identified per HIPAA Safe Harbor requirements before delivery to the Physical AI system. De-identification must remove or generalize all 18 categories of Protected Health Information as defined in 45 CFR 164.514(b)(2): names; geographic subdivisions smaller than a state; dates (except year) related to the individual; telephone numbers; fax numbers; email addresses; Social Security numbers; medical record numbers; health plan beneficiary numbers; account numbers; certificate/license numbers; vehicle identifiers and serial numbers; device identifiers and serial numbers; web universal resource locators; internet protocol addresses; biometric identifiers; full-face photographs and comparable images; and any other unique identifying number, characteristic, or code.

(b) \textit{Pseudonymization.} When re-linkable data is required for longitudinal tracking within the investigation, pseudonymization must use HMAC-SHA256 with site-specific salts. The pseudonymization key material must be stored separately from the pseudonymized data, with access restricted to authorized personnel at the clinical site. Physical AI systems must not have access to the pseudonymization key material.

(c) \textit{Deny-by-Default Authorization.} Physical AI systems must operate under deny-by-default role-based access control (RBAC) as enforced by the trialmcp-authz server. All data access requests must be explicitly authorized by policy before being granted. The six-actor permission matrix defined in the national MCP-PAI standard governs access: Robot Agent, Trial Coordinator, Data Monitor, Auditor, Sponsor, and CRO each have explicitly defined permissions. Explicit DENY rules take precedence over ALLOW rules in all cases.

(d) \textit{Hash-Chained Audit Trails.} All Physical AI system actions must be recorded in hash-chained immutable audit trails via the trialmcp-ledger server. Each audit record must contain the action performed, actor identity, timestamp (UTC), resource affected, outcome, and a SHA-256 hash computed over the canonical JSON representation of the record concatenated with the hash of the previous record. The genesis record uses a defined genesis hash. These audit trails satisfy the electronic records requirements of 21 CFR Part 11 for attributability, legibility, contemporaneousness, originality, and accuracy.

(e) \textit{Provenance Tracking.} The complete data lineage for every Physical AI procedure must be documented through DAG-based provenance tracking via the trialmcp-provenance server. Provenance records must trace the path from source clinical data (FHIR resources, DICOM studies) through all transformations, processing steps, and analytical operations to final clinical outcomes. Each provenance node must include a SHA-256 fingerprint of the data at that point in the lineage, enabling verification of data integrity throughout the processing chain.

(f) \textit{Federated Learning.} In multi-site investigations that employ federated learning for AI model training, the following data protections must be implemented: differential privacy mechanisms must be applied to model updates shared across sites, with documented epsilon and delta parameters; secure aggregation protocols (additive secret sharing, pairwise masking) must prevent any single site or aggregator from accessing another site's raw model updates; data harmonization across sites must be documented, including mappings between DICOM metadata formats, FHIR resource representations, ICD-10 to SNOMED CT coding, and LOINC tumor marker standardization; and privacy budget tracking must ensure that cumulative privacy expenditure across all federated learning rounds does not exceed the pre-specified privacy budget for the investigation.

\subsection{\S~50.34 Physical AI System Classification and Regulatory Pathways}

(a) \textit{FDA AI/ML SaMD Classification.} Physical AI systems that incorporate AI/ML components functioning as Software as a Medical Device (SaMD) are subject to FDA regulatory classification and oversight. The classification of the AI/ML component depends on the significance of the information provided by the software to the healthcare decision and the state of the healthcare situation or condition, consistent with the International Medical Device Regulators Forum (IMDRF) risk categorization framework and FDA guidance on AI/ML-enabled medical devices (January 2025).

(b) \textit{Regulatory Submission Pathways.} Depending on the risk classification of the Physical AI system, the following regulatory submission pathways may apply:

\begin{description}[leftmargin=1.0cm,style=sameline]
\item[(1)] \textit{510(k) Premarket Notification.} For Physical AI systems that are substantially equivalent to a legally marketed predicate device. This pathway may be appropriate for incremental AI/ML enhancements to existing robotic surgical platforms with established predicates.
\item[(2)] \textit{De Novo Classification.} For novel Physical AI systems of low to moderate risk that lack a suitable predicate. This pathway may be appropriate for novel AI-guided diagnostic or therapeutic robotic systems that represent new device types.
\item[(3)] \textit{Premarket Approval (PMA).} For high-risk Physical AI systems, including autonomous surgical robots and AI systems that independently make critical clinical decisions. PMA requires demonstration of safety and effectiveness through valid scientific evidence.
\item[(4)] \textit{Breakthrough Device Designation.} Physical AI systems that provide more effective treatment or diagnosis of life-threatening or irreversibly debilitating diseases may qualify for Breakthrough Device Designation, which provides intensive FDA interaction and priority review.
\end{description}

(c) \textit{Predetermined Change Control Plan (PCCP).} Adaptive Physical AI systems that are designed to learn and evolve based on clinical experience must maintain a documented PCCP that specifies: the types of modifications anticipated (e.g., AI model retraining, parameter updates, algorithm refinements); the methodology for implementing and validating each type of change; the performance metrics and acceptance criteria that modifications must satisfy before deployment; and the process for assessing the impact of modifications on safety and effectiveness.

(d) \textit{Applicable Safety Standards.} Physical AI systems used in clinical investigations must comply with applicable safety standards, including but not limited to:

\begin{description}[leftmargin=1.0cm,style=sameline]
\item[(1)] IEC 80601-2-77 for robotically assisted surgical equipment, which addresses basic safety and essential performance requirements for surgical robotic systems.
\item[(2)] ISO 14971 for application of risk management to medical devices, which provides the framework for identifying, evaluating, controlling, and monitoring risks associated with Physical AI systems throughout their lifecycle.
\item[(3)] ISO 13482 for safety requirements for personal care robots, which is applicable to rehabilitative exoskeletons, companion robots, and other Physical AI systems that interact directly with patients in non-surgical capacities.
\item[(4)] IEC 62304 for medical device software lifecycle processes, which governs the development, maintenance, and risk management of software components within Physical AI systems.
\item[(5)] ISO/TS 15066 for collaborative robot safety, which specifies safety requirements for cobots operating alongside clinical staff in shared workspaces.
\end{description}

\section{SUBPART D -- ADDITIONAL SAFEGUARDS FOR CHILDREN IN CLINICAL INVESTIGATIONS}

\textbf{Source:} 66 FR 20598, Apr. 24, 2001, unless otherwise noted.

\subsection{\S~50.50 IRB Duties}

In addition to other responsibilities assigned to IRBs under this subpart, each IRB must review proposed clinical investigations involving children as subjects covered by this subpart and approve only those clinical investigations that satisfy the criteria described in \S~50.51, \S~50.52, or \S~50.53, or meet the conditions described in \S~50.54.

\subsubsection{Physical AI Adaptation of \S~50.50}

When the clinical investigation involves Physical AI systems and children as subjects, the IRB must additionally evaluate whether the Physical AI system's safety profile, human oversight plan, and informed consent materials are appropriate for the pediatric population. The IRB must consider the child's age, developmental stage, and ability to understand and interact with robotic systems when evaluating the adequacy of assent procedures and Physical AI system configurations. For assistive robots used in pediatric oncology settings (companion robots, humanoid assistants), the IRB must evaluate the age-appropriateness of the robotic interaction protocols.

\subsection{\S~50.51 Clinical Investigations Not Involving Greater Than Minimal Risk}

An IRB may find that a clinical investigation involving children as subjects involves no greater than minimal risk, provided that adequate provisions are made for soliciting the assent of the children and the permission of their parents or guardians as set forth in \S~50.55.

\subsubsection{Physical AI Adaptation of \S~50.51}

When Physical AI systems are involved in clinical investigations with children that are proposed as involving no greater than minimal risk, the IRB must apply the Physical AI minimal risk guidelines specified in \S~50.22(f) of this part. Companion monitoring systems performing only non-contact observation and verbal interaction may qualify for minimal risk classification in pediatric settings if the system's interaction protocols have been specifically validated for the target age range. Robotic systems involving physical contact with children generally do not qualify for minimal risk classification.

\subsection{\S~50.52 Clinical Investigations Involving Greater Than Minimal Risk but Presenting the Prospect of Direct Benefit}

An IRB may find that a clinical investigation involving children as subjects involves greater than minimal risk but presents the prospect of direct benefit to the individual subjects, provided that:

(a) The risk is justified by the anticipated benefit to the subjects;

(b) The relation of the anticipated benefit to the risk is at least as favorable to the subjects as that presented by available alternative approaches; and

(c) Adequate provisions are made for soliciting the assent of the children and permission of their parents or guardians as set forth in \S~50.55.

\subsubsection{Physical AI Adaptation of \S~50.52}

When Physical AI systems are used in investigations involving children with the prospect of direct benefit, the risk-benefit assessment required under paragraph (a) must specifically account for Physical AI-specific risks (mechanical injury, AI prediction error, software malfunction) and Physical AI-specific benefits (enhanced surgical precision, reduced procedure duration, simulation-validated approaches). The comparison with alternative approaches required under paragraph (b) must include comparison with the same procedure performed without Physical AI assistance.

\subsection{\S~50.53 Clinical Investigations Involving Greater Than Minimal Risk and No Prospect of Direct Benefit}

An IRB may find that a clinical investigation involving children as subjects involves greater than minimal risk and no prospect of direct benefit to the individual subjects, provided that:

(a) The risk represents a minor increase over minimal risk;

(b) The intervention or procedure presents experiences to subjects that are reasonably commensurate with those inherent in their actual or expected medical, dental, psychological, social, or educational situations;

(c) The intervention or procedure is likely to yield generalizable knowledge about the subjects' disorder or condition that is of vital importance for the understanding or amelioration of the subjects' disorder or condition; and

(d) Adequate provisions are made for soliciting the assent of the children and permission of their parents or guardians as set forth in \S~50.55.

\subsubsection{Physical AI Adaptation of \S~50.53}

Physical AI systems involving physical contact with children should generally not be approved under this section unless the robotic interaction is limited to monitoring or rehabilitative functions that represent only a minor increase over minimal risk. The use of surgical robots, therapeutic positioning systems, or diagnostic needle-placement platforms on children without prospect of direct benefit to the individual child is not appropriate under this section.

\subsection{\S~50.54 Clinical Investigations Not Otherwise Approvable}

(a) The IRB finds that the clinical investigation presents a reasonable opportunity to further the understanding, prevention, or alleviation of a serious problem affecting the health or welfare of children; and

(b) The Commissioner of Food and Drugs, after consultation with a panel of experts in pertinent disciplines and following opportunity for public review and comment, determines either that the investigation satisfies the conditions of \S~50.51, \S~50.52, or \S~50.53, or that the investigation presents a reasonable opportunity to further the understanding, prevention, or alleviation of a serious problem affecting the health or welfare of children, the investigation will be conducted in accordance with sound ethical principles, and adequate provisions are made for soliciting the assent of children and the permission of their parents or guardians as set forth in \S~50.55.

\subsubsection{Physical AI Adaptation of \S~50.54}

When Physical AI systems are involved in investigations submitted to the Commissioner under this section, the panel of experts consulted must include at least one member with expertise in Physical AI systems, pediatric robotic surgery, or pediatric medical device safety. The public review and comment period must include materials explaining the Physical AI systems proposed for use with children in accessible, non-technical language.

\subsection{\S~50.55 Requirements for Permission by Parents or Guardians and for Assent by Children}

(a) In addition to the determinations required under other applicable sections of this subpart D, the IRB must determine that adequate provisions are made for soliciting the assent of the children when in the judgment of the IRB the children are capable of providing assent.

(b) In determining whether children are capable of providing assent, the IRB must take into account the ages, maturity, and psychological state of the children involved. This judgment may be made for all children to be involved in clinical investigations under a particular protocol, or for each child, as the IRB deems appropriate.

(c) The assent of the children is not a necessary condition for proceeding with the clinical investigation if the IRB determines that the capability of some or all of the children is so limited that they cannot reasonably be consulted, or that the intervention or procedure involved in the clinical investigation holds out a prospect of direct benefit that is important to the health or well-being of the children and is available only in the context of the clinical investigation.

(d) Even where the IRB determines that the subjects are capable of assenting, the IRB may still waive the assent requirement if it finds and documents that the clinical investigation involves no more than minimal risk, the waiver will not adversely affect the rights and welfare of the subjects, the investigation could not practicably be carried out without the waiver, and whenever appropriate, the subjects will be provided with additional pertinent information after participation.

(e) In addition to the determinations required under other applicable sections of this subpart D, the IRB must determine, in accordance with and to the extent that consent is required under part 50, that the permission of each child's parents or guardian is granted. Where parental permission is to be obtained, the IRB may find that the permission of one parent is sufficient for investigations conducted under \S~50.51 or \S~50.52. Where investigations are covered by \S~50.53 or \S~50.54 and permission is to be obtained from parents, both parents must give their permission unless one parent is deceased, unknown, incompetent, or not reasonably available, or when only one parent has legal responsibility for the care and custody of the child.

(f) Permission by parents or guardians must be documented in accordance with and to the extent required by \S~50.27.

(g) When the IRB determines that assent is required, it must also determine whether and how assent must be documented.

\subsubsection{Physical AI Adaptation of \S~50.55}

(h) When clinical investigations involve Physical AI systems and children as subjects, the assent process must include age-appropriate explanations of the robotic systems the child may encounter. For younger children who may not understand technical descriptions, visual aids including photographs, videos, or supervised demonstrations of the robotic system in a non-clinical setting should be provided. For older children and adolescents capable of understanding more detailed information, the assent materials should include a simplified version of the Physical AI-specific consent elements required under \S~50.25 of this part.

(i) Parental permission for Physical AI investigations must include all Physical AI-specific elements required in \S~50.25, including disclosure of the specific robot type, USL score, safety provisions, and the parent's right to request non-Physical AI alternatives for their child.

\subsection{\S~50.56 Wards}

(a) Children who are wards of the State or any other agency, institution, or entity can be included in clinical investigations approved under \S~50.53 or \S~50.54 only if such clinical investigations are related to their status as wards, or conducted in schools, camps, hospitals, institutions, or similar settings in which the majority of children involved as subjects are not wards.

(b) If the clinical investigation is approved under paragraph (a), the IRB must require appointment of an advocate for each child who is a ward. The advocate will serve in addition to any other individual acting on behalf of the child as guardian or in loco parentis. One individual may serve as advocate for more than one child. The advocate must be an individual who has the background and experience to act in, and agrees to act in, the best interest of the child for the duration of the child's participation in the clinical investigation. The advocate must not be associated in any way (except in the role as advocate or member of the IRB) with the clinical investigation, the investigator(s), or the guardian organization.

\subsubsection{Physical AI Adaptation of \S~50.56}

(c) When clinical investigations under this section involve Physical AI systems, the advocate appointed for each ward must be informed about the Physical AI systems used in the investigation, including the types of robotic systems, the degree of autonomous operation, the safety provisions, and the availability of non-Physical AI alternatives. The advocate must have access to the Physical AI-specific consent elements described in \S~50.25 and must be given adequate time and opportunity to understand the Physical AI aspects of the investigation before the ward's participation begins.

